\chapter{ОГЛЯД ЛІТЕРАТУРИ}

Промислові підприємства ставлять IoT для моніторингу обладнання та оптимізації виробництва. Однак виникає проблема: тисячі сенсорів генерують великі обсяги даних, які треба обробляти в реальному часі, знаходити аномалії та приймати рішення автоматично.

У статті \cite{alabadi2022industrial} автори описують Індустрію 4.0 --- це зміни у виробництві, торгівлі, транспорті та нафтовій галузі через Промисловий IoT. Це підвищує продуктивність і дає змогу обладнанню з різних локацій працювати разом.

Автори описали стан Індустрії 4.0 та IIoT з обмеженнями та досягненнями. Розглянули архітектуру, концепції та приклади застосування IIoT. Окрема увага була приділена технологіям, що підтримують Індустрію 4.0. Автори виділяють три напрями розвитку IIoT: глибинне навчання, периферійні обчислення та великі дані. На цій основі пропонують оновлену архітектуру майбутнього IIoT.

На даний момент великі обсяги даних у реальному часі вимагають системи, які не просто показують дані, а автоматично знаходять відхилення та потенційні збої через машинне навчання.

Є багато досліджень та комерційних рішень для промислових IoT-систем. Але інтеграція аналітики, візуалізації та обробки даних у одній платформі --- все ще проблема. Більшість систем фокусуються або на зборі телеметрії, або на аналізі даних. Деякі вимагають складного та дорогого налаштування. Треба універсальні та гнучкі рішення з зручним інтерфейсом та алгоритмами детекції аномалій.

\section{Архітектури для обробки потокових даних}

Є кілька підходів для обробки IoT даних у реальному часі. Кожен має плюси та мінуси.

\subsection{Lambda-архітектура}

У роботі \cite{schrabauer2021lambda} автори дослідили проблеми сучасних систем обробки IoT-даних. Показали, що централізовані рішення не підходять для великих додатків через зростання обсягів та швидкості даних.

Рішення --- Lambda-архітектура. Вона має три шари:

\begin{itemize}
    \item Пакетний шар (batch layer) --- зберігає історичні дані та формує точні результати
    \item Швидкий шар (speed layer) --- аналіз у реальному часі з мінімальною затримкою
    \item Обслуговуючий шар (serving layer) --- об'єднує результати обох шарів для відповіді на запити
\end{itemize}

Архітектура реалізована наступним чином. Швидкий шар працює на периферії через Apache Flink і Cassandra. Пакетний та обслуговуючий шари у хмарі через Hadoop, HDFS, MapReduce та HBase. Сенсори передають дані через Kafka на периферійні вузли.

Docker-контейнери та оркестрація через Ansible/Docker Swarm дають масштабованість та підтримку різної інфраструктури. Обчислення на периферії зменшують затримки порівняно з хмарою. Експерименти показують зменшення часу відгуку, особливо при затримках понад 50 мс або великих пакетах, що є однією з переваг.

Обмеження також присутні, адже складна підтримка через дублювання логіки між пакетним та швидким шарами. Кожна зміна вимагає оновлення обох компонентів. Це підвищує складність та ризик помилок у результатах.

\subsection{Kappa-архітектура}

На відміну від Lambda, Kappa-архітектура \cite{nkamla2022optimized} використовує єдиний потоковий конвеєр, що спрощує підтримку системи, але вимагає повторної обробки історичних даних при зміні логіки.

Переваги:
\begin{itemize}
    \item Єдиний шлях коду
    \item Простіша підтримка та розгортання
    \item Менша операційна складність
\end{itemize}

Обмеження:
\begin{itemize}
    \item Повторна обробка всіх даних при зміні логіки
    \item Вищі вимоги до сховища для повторної обробки
    \item Складність пакетного аналізу на історичних даних
\end{itemize}

\subsection{Архітектури для ML-систем}

У статті \cite{vajpayee2024building} розглядається роль масштабованих архітектур даних у машинному навчанні. Обсяги даних ростуть, моделі ускладнюються. Автори описують основні компоненти таких архітектур: збір, зберігання, обробка даних та розгортання моделей. Також ключові шаблони --- Lambda, Kappa та мікросервіси. Вони обговорюють технології та інструменти, необхідні для реалізації масштабованих ML-систем, та наголошують на інтеграції машинного навчання з процесами обробки даних.

Як приклад --- кейс предиктивного обслуговування у виробництві. Він показує практичну користь архітектури: скорочення часу простою обладнання та економія витрат. Масштабовані архітектури дозволяють ефективно управляти великими потоками даних та підвищувати продуктивність ML-рішень.

Сучасні гібридні підходи поєднують потокову обробку з матеріалізованими представленнями (ksqlDB, Apache Flink). Це дозволяє аналітику в реальному часі без складності Lambda.

\section{Промисловий IoT моніторинг}

Готові платформи спрощують інтеграцію та скорочують час розробки \cite{ferencz2024cloud}.

\subsection{Існуючі платформи}

AWS IoT --- управління пристроями, обробка даних у реальному часі, інтеграція з AI. Плюси: готові інструменти, масштабування, надійність. Мінуси: висока вартість, прив'язка до провайдера.

Azure IoT --- IoT Hub, потокова аналітика, інтеграція з машинним навчанням. Плюси: інтеграція з екосистемою Microsoft, корпоративна підтримка. Мінуси: складне налаштування, високі витрати при масштабуванні.

ThingsBoard --- відкритий код з підтримкою MQTT, HTTP, CoAP, панелями керування та правилами обробки. Плюси: гнучкість, без ліцензійних витрат. Мінуси: треба своя команда для адміністрування.

\subsection{Порівняння функціональності}

Комерційні платформи дають автоматичне масштабування та вбудовану безпеку, але дорогі та створюють прив'язку до постачальника. Відкритий код дає повний контроль без ліцензійних витрат, але вимагає власної інфраструктури та DevOps-команди.

\subsection{Виклики та обмеження}

Основні виклики: управління великими обсягами даних від тисяч сенсорів, інтеграція різних пристроїв з різними протоколами (Modbus, OPC UA, MQTT), безпека промислових систем та аналітика у реальному часі з виявленням аномалій. Гібридний підхід (комбінація хмарних і платформ з відкритим кодом) дає баланс між функціональністю, витратами та гнучкістю.


\section{Алгоритми виявлення аномалій}

\subsection{Підходи без вчителя для часових рядів}

У роботі \cite{audibert2020usad} представлено USAD (виявлення аномалій без вчителя) --- алгоритм для виявлення аномалій у мультиваріантних часових рядах. Він базується на автоенкодерах та навчанні з елементами GAN.

 Не потребує мічених аномалій для навчання, що критично для промислових систем, де аномалії рідкісні та важко задокументовані. Стабільне та швидке навчання навіть при різних налаштуваннях параметрів. Одна модель може генерувати кілька рівнів чутливості для різних сценаріїв моніторингу, що є очевидною перевагою даного підходу.

Як обмеження можна виділити, що це вимагає періодів без аномалій для навчання. Потрібні ресурси для обробки мультиваріантних часових рядів у реальному часі. Треба налаштовувати чутливість, щоб уникнути перевантаження хибними тривогами.

\subsection{Isolation Forest для виявлення аномалій}

На відміну від більшості модельних підходів, які будують профіль нормальних екземплярів та ідентифікують аномалії як відхилення від цього профілю, Isolation Forest (iForest) \cite{liu2008isolation} пропонує фундаментально інший метод, який явно ізолює аномалії замість профілювання нормальних точок. Концепція ізоляції базується на спостереженні, що аномалії --- це рідкісні та відмінні спостереження, які легше ізолювати в просторі ознак через меншу кількість умов, необхідних для їх відокремлення.

 Використання принципу ізоляції дозволяє iForest застосовувати підвибірку в обсязі, недоступному для існуючих методів, створюючи алгоритм з лінійною часовою складністю, низькою константою та низькими вимогами до пам'яті. Алгоритм ефективно працює у високовимірних задачах з великою кількістю нерелевантних атрибутів, а також у ситуаціях, де навчальна вибірка не містить жодних аномалій.

Емпіричні оцінки показують, що iForest демонструє кращу продуктивність порівняно з ORCA (метод на основі відстані з майже лінійною складністю), LOF та Random Forest за метриками AUC та часу обробки, особливо на великих наборах даних.

\subsection{Підходи з вчителем для IoT-систем з незбалансованими даними}

Дослідження \cite{doghramachi2023internet} присвячене проблемі виявлення атак в промислових IoT середовищах, де більшість доступних наборів даних страждають від незбалансованості класів. Запропонована методологія включає попередню обробку даних (очищення, кодування, нормалізацію), застосування методу синтетичної надвибірки меншості (SMOTE) для вирішення проблеми незбалансованості класів, та зменшення розмірності через кореляційний аналіз.

Використані алгоритми:
\begin{itemize}
\item Логістична регресія --- базова модель для класифікації
\item Багатошаровий персептрон --- нейронна мережа для складних патернів
\item Дерево рішень --- інтерпретована деревоподібна структура
\item Випадковий ліс (Random Forest) --- ансамбль дерев для стабільності
\item XGBoost --- градієнтний бустинг для максимальної точності
\end{itemize}

Результати на наборі даних IoTID20: XGBoost з інтеграцією SMOTE досягає видатної точності виявлення атак: 0.99 для бінарної класифікації, 0.99 для мультикласової класифікації, та 0.81 для множинної суб-класифікації. Ці результати демонструють значне покращення порівняно з існуючими фреймворками систем виявлення вторгнень для незбалансованих IoT наборів даних.

\subsection{Порівняння підходів для IoT}

Аналіз літератури виявляє взаємодоповнюваність підходів без вчителя та з вчителем для IoT-систем. Методи без вчителя (USAD, Isolation Forest) ефективні для промислового моніторингу, де аномалії рідкісні та важко задокументовані --- Isolation Forest забезпечує лінійну часову складність, низькі вимоги до пам'яті та здатність працювати без мічених аномалій у навчальній вибірці, що критично для реальних промислових систем. Методи з вчителем (Random Forest, XGBoost) демонструють вищу точність при наявності збалансованих навчальних даних: XGBoost з SMOTE досягає 0.99 точності для класифікації відомих типів атак в промисловому IoT, перевершуючи традиційні підходи систем виявлення вторгнень. Для реальних IoT-систем вибір залежить від наявності мічених даних та характеру задачі: без вчителя (Isolation Forest) для виявлення невідомих аномалій без потреби в навчальних даних, з вчителем (XGBoost) для високоточної класифікації відомих загроз при наявності репрезентативних наборів даних.


\section{Практики MLOps}

MLOps (операції машинного навчання) --- підхід для керування, автоматизації та операціоналізації процесів розробки, розгортання та підтримки ML-моделей \cite{hanchuk2025implementing}. Ці практики забезпечують стабільність, відтворюваність та ефективність моделей у виробничому середовищі.

\subsection{Версіонування та реєстр моделей}

В даній статті, автори поставили собі за мету -- забезпечити відстежуваність змін моделей та даних у часі, підтримувати контроль якості та відтворюваність результатів.

Кожна версія моделі реєструється у реєстрі моделей (MLflow Model Registry, Kubeflow Pipelines). Включає метадані: версія моделі, параметри навчання, метрики продуктивності, артефакти (ваги, конфігурації конвеєрів). Підтримує відкати --- повернення до попередньої стабільної версії при виникненні проблем.

Це дозволяє відтворювати результати експериментів, інтегрувати моделі у конвеєр CI/CD та керувати повним життєвим циклом моделі.

\subsection{Безперервне навчання}

Забезпечення постійного оновлення моделі у відповідь на нові дані та зміни у середовищі.

Автоматизовані конвеєри для регулярного тренування моделей --- заплановане перенавчання або перенавчання за тригером при надходженні нових даних. Використання CI/CD для ML, де нові версії моделі автоматично тренуються та тестуються перед розгортанням. Це критично для моделей, що залежать від динамічних даних або швидко змінюваних умов.

Інструменти: MLflow, Kubeflow, Polyaxon, Seldon Core.

\subsection{Моніторинг та виявлення дрейфу}

Контроль продуктивності моделей у виробничому середовищі та виявлення деградації або змін даних.

Моніторинг ключових метрик: точність, повнота, затримка, використання ресурсів. Виявлення дрейфу даних --- зміни у розподілі вхідних даних, що впливають на якість прогнозів. Виявлення дрейфу концепції --- зміни у зв'язку між ознаками та цільовою змінною, що знижують ефективність моделі. Автоматичне сповіщення та тригери для перетренування при виявленні дрейфу.

Інструменти: MLflow (моніторинг моделей), Prometheus + Grafana, Seldon Analytics, Evidently AI.

\subsection{Висновок}

Практики MLOps дозволяють підвищити стабільність та відтворюваність моделей, автоматизувати оновлення та підтримку у виробничому середовищі, швидко реагувати на зміни у даних та середовищі. Це критично для ефективності ML-рішень у промисловому IoT моніторингу.


\section{Мультитенантна архітектура}

Мультитенантність дозволяє одному екземпляру програми обслуговувати кількох користувачів (тенантів) із логічною ізоляцією даних \cite{kumar2020multi}. Це забезпечує масштабованість, ефективність і економію ресурсів у хмарних середовищах. Використовується у SaaS-платформах (ERP, CRM, HR) та IoT-бекендах.

\subsection{Моделі ізоляції даних}

База даних на тенанта. Кожен тенант має власну базу даних. Висока ізоляція даних, але висока вартість через окремі екземпляри БД. Середня масштабованість --- складно керувати сотнями баз. Підходить для регульованих галузей (фінанси, медицина), де критична повна сегрегація даних.

Схема на тенанта. Одна база даних, але кожен тенант має окрему схему. Середня ізоляція --- дані розділені логічно, але на одному сервері. Висока масштабованість --- простіше керувати, ніж окремими БД. Середня вартість. Підходить для корпоративних SaaS з помірними вимогами до ізоляції.

Таблиця на тенанта (спільна база даних). Усі тенанти в одній базі та схемі, розділення через ідентифікатор тенанта у кожній таблиці. Низька ізоляція --- ризик помилок у запитах може призвести до витоку даних. Дуже висока масштабованість --- один екземпляр БД обслуговує тисячі тенантів. Низька вартість. Підходить для великих SaaS-платформ та маркетплейсів з менш чутливими даними.

\subsection{Міркування щодо безпеки}

Ізоляція даних. Основний виклик --- запобігання випадковому чи несанкціонованому доступу між тенантами. Потрібна ретельна валідація ідентифікатора тенанта у кожному запиті. Помилка у фільтрації може призвести до витоку даних іншого тенанта.

Контроль доступу. Складність із керуванням доступом на основі ролей (RBAC) або атрибутів (ABAC) для великої кількості користувачів і ролей. Треба керувати правами на рівні тенанта, користувача та ресурсу одночасно.

Масштабування та навантаження. Один тенант може створювати надмірне навантаження, що впливає на інших. Потрібні механізми обмеження ресурсів та обмеження швидкості запитів.

Регуляторні вимоги. GDPR, HIPAA, SOC 2 вимагають сегрегації даних та повного аудиту доступу. Для регульованих галузей краще використовувати моделі база даних на тенанта або схема на тенанта.

Вибір моделі залежить від балансу між безпекою, вартістю та масштабованістю конкретного випадку використання.
